\chapter*{Zusammenfassung}
% \thispagestyle{empty}
Die vorliegende Arbeit befasst sich mit der modellbasierten Funktionsentwicklung für Akkugeräte. Ziel ist es, zukünftig Funktionen für Akkugeräte modellbasiert in Matlab zu entwickeln und diese über die Autocodegenerierung in die Firmware zu integrieren. Im Rahmen dieser Arbeit werden erste Funktionen in die Simulationsumgebung übertragen und der Grundstein für die modellbasierte Funktionsentwicklung gelegt. Zunächst wird eine Modellstruktur sowie ein Ablauf für die Code-Generierung erarbeitet. Anschließend werden erste Maschinenschutzfunktionen in Matlab modelliert, in der Simulationsumgebung getestet und unter Berücksichtigung der Anforderungen aus dem Lastenheft verifiziert. In weiteren Systemtests werden an einer realen akkubetriebenen Motorsäge die Funktionen validiert.\\
Die abschließende Speicher- und Laufzeitanalyse zeigt, dass die modellbasierten Funktionen eine geringe Erhöhung des Ressourcenbedarfs im Vergleich zu den manuell codierten Funktionen verursachen. Dies ist vor allem auf die Schnittstelle zwischen dem Autocode und der Basissoftware zurückzuführen. Die Ergebnisse dieser Arbeit bieten eine solide Grundlage für den Übergang von der manuellen Programmierung hin zur modellbasierten Entwicklung, hinsichtlich nicht-sicherheitsrelevanter Funktionen.

%Das muss vor das Inhaltsverzeichnis und die ganzen anderen Verzeichnisse. 
%Die vorliegende Arbeit beschäftigt sich mit der Entwicklung eines bionisch inspirierten Multigelenk-Roboters für die Medizintechnik. Der Roboter ist "`schlangenförmig"' aufgebaut und kann bei eingeschränkten Platzverhältnissen eingesetzt werden. Aufgrund der bei medizinischen Anwendungen gegebenen räumlichen Begrenzungen weißt der Roboter eine hohe Beweglichkeit auf.     
%
%Die vorliegende Arbeit gliedert sich in ... Abschnitte, anhand derer die Entwicklung und Inbetriebnahme des Roboters beschrieben werden.

%%%%%%%%%%%%%%%%
% Die vorliegende Arbeit beschäftigt sich mit einer ersten Konzeptionierung und Entwicklung eines bionisch inspirierten Multigelenk-Roboters für die Medizintechnik. Der Roboter soll schlangenförmig aufgebaut und bei eingeschränkten Platzverhältnissen eingesetzt werden können. Die Aktuierung erfolgt mittels Formgedächtnisdrähten. Dafür werden zunächst im Rahmen einer umfassenden Literaturrecherche bereits bestehende Multigelenk-Roboter hinsichtlich ihrer Gelenkmechanismen untersucht und mögliche Anwendungsgebiete eines solchen Roboters aufgezeigt. Zudem werden grundlegende Überlegungen bezüglich einer Realisierung eines Multigelenk-Roboters aufgestellt. Diese Überlegungen werden mit unterschiedlichen Bewertungs- und Auswahlverfahren zu einem geeigneten Realisierungskonzepts kombiniert.\\
% Durch die Entwicklung und Inbetriebnahme eines ersten Funktionsmodells wird der gewählte Ansteuermechanismus, der die Bewegung des Roboterarms ausführt, getestet. Aus diesen Versuchen werden wichtige Erkenntnisse für das ausgewählte Realisierungskonzept gewonnen. Das Grundprinzip des Ansteuermechanismus besteht darin, die verwendbare ausgeführte Bewegung der Formgedächtnisdrähte zu vergrößern.\\
% Nach erfolgreicher Implementierung des Ansteuermechanismus im Funktionsmodells erfolgt abschließend eine Machbarkeintsanalyse, die das Erreichen der Zielbaugröße veranschaulicht. Durch das erarbeitete Konzept und der in der Theorie erfolgreichen Erreichbarkeit der Zielbaugröße, ist ein erster Grundstein für die Entwicklung eines neuartigen Multigelenk-Roboters gelegt.
%%%%%%%%%%%%%%%%%%%

% Aufgrund der räumlichen Begrenzungen in medizinischen Anwendungen weist der Roboter eine hohe Beweglichkeit auf.\\
%Die Arbeit gliedert sich in sechs Abschnitte, die die Entwicklung und Inbetriebnahme des Roboters beschreiben. Zunächst werden das Problem und die Anforderungen an das System erläutert. Anschließend wird eine umfassende Literaturrecherche durchgeführt, die in drei Hauptbereiche unterteilt ist: Bionik, (Anwendungen??) und Einsatzgebiete.
%
%Das darauf folgende Kapitel behandelt die Konzipierung des Multigelenk-Roboters. Es umfasst die grundlegenden Überlegungen und Möglichkeiten bis hin zur Bewertung und Auswahl eines geeigneten Realisierungskonzepts. Die Sensorik für den Roboter wird in einem eigenen Kapitel behandelt. Darin werden verschiedene Sensortechnologien für eine Implementierung im Zielkonzept sowie für ein Funktionsmodell vorgestellt und bewertet.
%
%Anschließend wird ein Funktionsmodell entwickel t, um ein neues Ansteuerungskonzept mittels SMA-Draht zu realisieren und dessen mögliche Umsetzung zu verifizieren. Dieses Modell wird schrittweise verbessert, bis die erwarteten und gewünschten Ergebnisse erreicht sind. Dabei wird auch ein Simulationsmodell der SMA-Drähte verwendet, um mögliche Fehler frühzeitig zu erkennen und Erwartungen zu formulieren.
%
%Zum Abschluss der Arbeit wird überprüft, ob das Antriebskonzept des Funktionsmodells in der Zielgröße umgesetzt werden kann.
